\documentclass[11pt,a4paper]{article}

\usepackage[T1]{fontenc}
\usepackage[utf8]{inputenc}
\usepackage[english]{babel}

\usepackage{amsmath,amssymb,amsthm}
\usepackage{graphicx}
\usepackage{hyperref}
\usepackage{geometry}
\usepackage{booktabs}

\geometry{margin=2.5cm}

\title{
	A Dirac-Type Operator from the Logarithmic Geometry of Primes
}

\author{
	Thomas Hoffbauer
}

\date{\today}

\begin{document}
	
	\maketitle
	
	\begin{abstract}
		
		We propose a Dirac-type operator constructed from the logarithmic geometry of prime numbers.
		Numerical experiments suggest that the spectrum of this operator
		is strongly correlated with the ordinates of the nontrivial zeros
		of the Riemann zeta function.
		
		This observation motivates a spectral framework in the spirit of the
		Hilbert--Pólya program.
		
	\end{abstract}
	
	\tableofcontents
	
	\section{Introduction}
	
	The nontrivial zeros of the Riemann zeta function
	
	\[
	\zeta(s)
	\]
	
	play a fundamental role in number theory.
	
	The Riemann hypothesis states
	
	\[
	\Re(s)=\frac12
	\]
	
	for all nontrivial zeros.
	
	One possible explanation is the Hilbert--Pólya idea:
	the zeros arise as eigenvalues of a self-adjoint operator.
	
	This article proposes a simple operator constructed from
	the logarithmic geometry of primes.
	
	---
	
	\section{Logarithmic Geometry of the Primes}
	
	Let
	
	\[
	p_1,p_2,p_3,\dots
	\]
	
	be the sequence of primes.
	
	Define the logarithmic coordinate
	
	\[
	C=\log p.
	\]
	
	This transformation linearizes multiplicative structures.
	
	The prime set becomes
	
	\[
	\mathcal P=\{\log p_1,\log p_2,\dots\}.
	\]
	
	---
	
	\section{Prime Graph}
	
	We construct a weighted graph
	
	\[
	G=(V,E)
	\]
	
	with vertices
	
	\[
	V=\{p_1,\dots,p_N\}.
	\]
	
	The edge weights are defined by
	
	\[
	w(p_i,p_j)=
	\frac{\log(p_i p_j)}{\sqrt{p_i p_j}}.
	\]
	
	This combines
	
	\begin{itemize}
		\item logarithmic scaling
		\item decay with distance
	\end{itemize}
	
	---
	
	\section{The BM Dirac Operator}
	
	Let
	
	\[
	Q_{ij}=w(p_i,p_j).
	\]
	
	Define
	
	\[
	D_{BM}=
	\begin{pmatrix}
		0 & Q\\
		Q^T & 0
	\end{pmatrix}.
	\]
	
	Properties:
	
	\begin{itemize}
		\item self-adjoint
		\item real spectrum
	\end{itemize}
	
	\[
	D_{BM}\psi_n=\lambda_n\psi_n.
	\]
	
	---
	
	\section{Numerical Experiments}
	
	\subsection{Spectrum}
	
	Compare
	
	\[
	\lambda_n
	\]
	
	with ordinates
	
	\[
	\gamma_n
	\]
	
	of the Riemann zeros.
	
	---
	
	\subsection{Spacing Statistics}
	
	Spacing
	
	\[
	s_n=\lambda_{n+1}-\lambda_n
	\]
	
	
	---
	
	\subsection{Unfolded Spectrum}
	
	Define unfolded eigenvalues
	
	\[
	\tilde{\lambda}_n
	\]
	
	by removing smooth density.
	
	Histogram of
	
	\[
	s_n=\tilde{\lambda}_{n+1}-\tilde{\lambda}_n
	\]
	
	is computed.
	
	---
	
	\subsection{Trace Formula Test}
	
	Compare
	
	\[
	\mathrm{Tr}\,f(D_{BM})
	\]
	
	with prime sums
	
	\[
	\sum_p \log(p)f(\log p).
	\]
	
	
	\section{Clifford Structure}
	
	Introduce the four-vector
	
	\[
	(A,C,e,b).
	\]
	
	Define
	
	\[
	X=A\gamma_0+C\gamma_1+b\gamma_2+e\gamma_3.
	\]
	
	This corresponds to the Clifford algebra
	
	\[
	Cl_{1,3}.
	\]
	
	---
	
	\section{Quaternion Interpretation}
	
	Define quaternion
	
	\[
	q=A+C\,i+b\,j+e\,k.
	\]
	
	Quaternion rotations correspond to SU(2) symmetry.
	
	---
	
	\section{Relation to Spectral Geometry}
	
	The BM construction resembles spectral triple ideas introduced by Alain Connes.
	
	The structure consists of
	
	\begin{itemize}
		\item algebra generated by primes
		\item Hilbert space of prime functions
		\item Dirac-type operator
	\end{itemize}
	
	---
	
	\section{Extended Geometric Interpretation}
	
	Possible extensions include
	
	\begin{itemize}
		\item quaternion structures
		\item octonionic extensions
		\item exceptional symmetry groups.
	\end{itemize}
	
	These ideas remain speculative.
	
	---
	
	\section{Conclusion}
	
	We proposed a simple operator derived from
	the logarithmic geometry of primes.
	
	Numerical experiments suggest a strong relation between
	its spectrum and the ordinates of the Riemann zeros.
	
	Further investigation of this operator may reveal
	deeper connections between primes, spectral theory,
	and geometry.
	
	---
	
	\begin{thebibliography}{9}
		
		\bibitem{riemann}
		B. Riemann,
		On the number of primes less than a given magnitude,
		1859.
		
		\bibitem{montgomery}
		H. Montgomery,
		Pair correlation of zeros of the zeta function.
		
		\bibitem{odlyzko}
		A. Odlyzko,
		Tables of Riemann zeta zeros.
		
		\bibitem{connes}
		A. Connes,
		Trace formula in noncommutative geometry.
		
	\end{thebibliography}
	
\end{document}